\subsection{Analyse Comparative des Solutions Industrielles}

\subsubsection{Agrona et la gestion de la mémoire "Off-Heap"}

\paragraph{Optimisation des structures de données via Agrona :}

Pour pallier l'inefficience des collections de la JDK standard, nous intégrons la bibliothèque \textit{Agrona}. Les collections \texttt{java.util} souffrent de deux tares rédhibitoires pour la très basse latence : l'impossibilité de stocker des types primitifs (induisant un \textit{autoboxing} massif et une pression sur le \textit{Garbage Collector}) et la dispersion en mémoire des données entraînant un phénomène de \textit{pointer chasing}.

Chaque objet Java possède un en-tête (\textit{Object Header}) de 12 à 16 octets, ajoutant une surcharge informationnelle et force le processeur à suivre des références vers des zones mémoires non contiguës. Les collections d'Agrona reposent sur des structures de données stockées dans des tableaux de primitifs contigus. Cette architecture favorise la localité spatiale : en plaçant les clés et les valeurs (dans le cadre des \textit{maps} et des \textit{sets}) de manière adjacente, Agrona maximise la probabilité qu'un couple de données réside dans une même \textit{cache line} de 64 octets. Ce design permet au \textbf{Hardware Prefetcher} du CPU d'anticiper les accès mémoire et de charger les données dans le cache L1 avant même leur sollicitation, masquant ainsi la latence de la DRAM.

\paragraph{Communication inter-threads et Ring Buffers :}

Pour la communication \textit{inter-threads} nous utilisons le \texttt{OneToOneRingBuffer} d'Agrona. Contrairement aux files d'attente bloquantes de la JDK, cette structure est \textit{lock-free} et \textit{wait-free} pour le producteur. Conformément au \textit{pattern Disruptor}, nous aurons un producteur unique ayant les accès en écriture et des consommateurs qui feront des appels au \texttt{OneToOneRingBuffer} pour les accès en lecture. 

\subsubsection{Chronicle Software : Persistance native et Communication Inter-Processus (\textit{IPC})}

Pour un \textit{matching engine}, la latence réseau entre la réception d'un ordre et son enregistrement est un facteur de compétitivité critique. L'approche de \textit{Chronicle Software} est de fusionner la communication entre les processus et la sauvegarde des données en une seule opération ultra-rapide basée sur la mémoire partagée.

\paragraph{Le fichier comme vecteur de communication et de persistance :}

Dans une architecture classique, un ordre doit être envoyé via un \textit{socket} réseau puis sauvegardé dans une base de données (disque). Chronicle supprime cette double étape en utilisant les fichiers mappés en mémoire (\textit{Memory Mapped Files}). Le fichier sur disque est projeté directement dans la RAM. Écrire un message pour le transmettre à un autre processus revient donc à l'écrire directement dans le fichier en RAM qui sera directement visible.

L'extrait suivant illustre comment l'écriture devient à la fois un acte de communication et un acte d'archivage :

\begin{lstlisting}[language=Java, caption=Ecriture Off-Heap persistante avec Chronicle-Core]
// 1. Ouverture du fichier journal (qui servira de canal de communication)
FileChannel fc = new CleaningRandomAccessFile(fileName, "rw").getChannel();

// 2. Mappage : le fichier devient une zone de RAM partagee entre les processus
long address = OS.map(fc, MapMode.READ_WRITE, 0, 64 << 10);

// 3. Persistance instantanee : l'ordre est ecrit en RAM, mais deja lie au disque
OS.memory().writeLong(address + 1024L, 0x1234567890ABCDEFL);

// 4. L'ordre est persiste et disponible pour les autres processus sans latence reseau
\end{lstlisting}

Comme la \textit{Gateway} et le \textit{Matching Engine} regardent le même fichier projeté en RAM, l'ordre n'est jamais déplacé. Sa présence dans le fichier suffit à le rendre "public" pour les autres composants. Chaque octet écrit est virtuellement déjà sur le disque. En cas de coupure de courant, le noyau Linux garantit que les données présentes dans cette RAM (le \textit{Page Cache}) sont les dernières enregistrées, permettant une reconstruction immédiate du carnet d'ordres au redémarrage.

\paragraph{S'affranchir du \textit{Garbage Collector} et du disque :}

La force de Chronicle est de rendre la persistance "invisible" pour le processeur. Le \textit{matching engine} ne subit jamais de ralentissement dû aux accès disques ou aux nettoyages de mémoire. L'écriture se fait à la vitesse de la RAM (nanosecondes). C'est le système d'exploitation qui, de manière asynchrone et en arrière-plan, synchronise physiquement les données sur le support disque. Le thread critique du moteur de matching ne bloque jamais pour attendre le disque. En utilisant la mémoire hors-tas (\textit{Off-Heap}), Chronicle ne crée aucun objet Java pour transporter l'ordre. On évite ainsi l'indéterminisme du \textit{Garbage Collector}.

\begin{lstlisting}[language=Java, caption=Gestion deterministe et atomique]
Memory memory = OS.memory();
long address = memory.allocate(1024);
try {
// Modification atomique (CAS) : coordination sans verrou (lock-free)
boolean swapped = memory.compareAndSwapInt(address, 1, 2);
} finally {
// Liberation manuelle : aucun nettoyage automatique imprevisible
memory.freeMemory(address, 1024);
}
\end{lstlisting}

L'usage de \texttt{compareAndSwapInt} (CAS) permet à plusieurs threads (par exemple le moteur de matching et le thread de reporting) de lire et modifier le journal d'ordres simultanément sans jamais se bloquer mutuellement. Le pipeline du processeur reste saturé d'instructions utiles, maximisant le débit d'ordres tout en garantissant que chaque transaction est immortalisée sur disque sans aucun sacrifice de vitesse.